\documentclass{article}
\usepackage{amsmath,amssymb,amsthm}
\usepackage{algorithm}
\usepackage{algorithmic}
\usepackage{hyperref}
\usepackage{enumitem}
\usepackage{geometry}
\geometry{margin=1in}
\newtheorem{theorem}{Theorem}
\newtheorem{lemma}{Lemma}
\newtheorem{assumption}{Assumption}
\newcommand{\EE}{\mathbb{E}}
\newcommand{\norm}[1]{\left\|#1\right\|}
\newcommand{\ip}[2]{\langle #1,#2\rangle}
\newcommand{\prox}{\operatorname{prox}}
\begin{document}

\section*{Algorithm Name}
\textbf{Extra‑Gradient Method for Convex–Concave Saddle‑Point Problems}  

We consider the problem  

\[
\min_{x\in\mathcal X}\;\max_{y\in\mathcal Y} \; L(x,y),
\]

where $L:\mathcal X\times\mathcal Y\rightarrow\mathbb R$ is convex in $x$,
concave in $y$, and continuously differentiable.  

The extra‑gradient method (also called the Korpelevich method) uses two
gradient evaluations per iteration.

\section*{Mathematical Formulation}
Let $z:=(x,y)$ and define the monotone operator  

\[
F(z)\;:=\;\begin{pmatrix}
\nabla_{x}L(x,y)\\[2mm]
-\nabla_{y}L(x,y)
\end{pmatrix}.
\]

Given a stepsize $\eta>0$, the iterates $\{z_k\}_{k\ge 0}$ are generated as

\[
\boxed{
\begin{aligned}
z_{k+\frac12}&:=z_k-\eta\,F(z_k) \qquad\text{(extragradient step)}\\[1mm]
z_{k+1}&:=z_k-\eta\,F(z_{k+\frac12}) \qquad\text{(correction step)} .
\end{aligned}}
\tag{EG}
\]

In component form,

\[
\begin{aligned}
x_{k+\frac12}&=x_k-\eta\,\nabla_xL(x_k,y_k),\\
y_{k+\frac12}&=y_k+\eta\,\nabla_yL(x_k,y_k),\\[1mm]
x_{k+1}&=x_k-\eta\,\nabla_xL(x_{k+\frac12},y_{k+\frac12}),\\
y_{k+1}&=y_k+\eta\,\nabla_yL(x_{k+\frac12},y_{k+\frac12}).
\end{aligned}
\tag{1}
\]

\section*{Pseudocode}
\begin{algorithm}[H]
\caption{Extra‑Gradient for $\min_x\max_y L(x,y)$}
\begin{algorithmic}[1]
\STATE \textbf{Input:} stepsize $\eta>0$, initial point $z_0=(x_0,y_0)$,
maximum iterations $T$.
\FOR{$k=0,\dots,T-1$}
    \STATE Compute the first gradient:
    \[
    g_k:=F(z_k)=\begin{pmatrix}\nabla_xL(x_k,y_k)\\ -\nabla_yL(x_k,y_k)\end{pmatrix}.
    \]
    \STATE Extragradient step:
    \[
    z_{k+\frac12}=z_k-\eta g_k .
    \]
    \STATE Compute the second gradient at the extragradient point:
    \[
    \tilde g_k:=F(z_{k+\frac12}) .
    \]
    \STATE Correction step:
    \[
    z_{k+1}=z_k-\eta\tilde g_k .
    \]
\ENDFOR
\STATE \textbf{Output:} $\displaystyle \bar z_T:=\frac1T\sum_{k=0}^{T-1}z_k$ (ergodic average).
\end{algorithmic}
\end{algorithm}

\section*{Dry Run Example}
Although the method is intended for a saddle‑point problem, the same
two‑step scheme can be applied to a simple convex function.
Consider  

\[
f(w)=(w-3)^2 ,\qquad w\in\mathbb R,
\]

with stepsize $\eta=0.1$ and initial point $w_0=0$.
The gradient is $\nabla f(w)=2(w-3)$.

\begin{center}
\begin{tabular}{c|c|c|c}
Iteration $k$ & $w_{k+\frac12}=w_k-\eta\nabla f(w_k)$ &
$w_{k+1}=w_k-\eta\nabla f(w_{k+\frac12})$ &
$f(w_{k+1})$ \\ \hline
0 & $0-0.1\cdot(-6)=0.6$ &
$0-0.1\cdot 2(0.6-3)=0-0.1\cdot(-4.8)=0.48$ &
$(0.48-3)^2=6.3504$ \\[2mm]
1 & $0.48-0.1\cdot2(0.48-3)=0.48-0.1\cdot(-5.04)=0.984$ &
$0.48-0.1\cdot2(0.984-3)=0.48-0.1\cdot(-4.032)=0.8832$ &
$(0.8832-3)^2=4.4615$ \\[2mm]
2 & $0.8832-0.1\cdot2(0.8832-3)=0.8832-0.1\cdot(-4.2336)=1.30656$ &
$0.8832-0.1\cdot2(1.30656-3)=0.8832-0.1\cdot(-3.38688)=1.221888$ &
$(1.221888-3)^2=3.1629$
\end{tabular}
\end{center}
The objective value decreases from $9$ (at $w_0=0$) to $\approx3.16$ after three
iterations, illustrating the effect of the extra‑gradient correction.

\newpage
\section*{Assumptions}
\begin{assumption}[Convex–concave structure]\label{as:convex}
For all $y\in\mathcal Y$, $x\mapsto L(x,y)$ is convex; for all $x\in\mathcal X$,
$y\mapsto L(x,y)$ is concave.
\end{assumption}
\begin{assumption}[Smoothness]\label{as:smooth}
There exists $L>0$ such that for all $z,z'\in\mathcal Z:=\mathcal X\times\mathcal Y$,
\[
\norm{F(z)-F(z')}\le L\,\norm{z-z'} .
\tag{2}
\]
Thus $F$ is $L$‑Lipschitz (i.e. $L$‑smooth).
\end{assumption}
\begin{assumption}[Bounded domain]\label{as:bounded}
The feasible sets $\mathcal X,\mathcal Y$ are convex and compact. Define
\[
D:=\max_{z\in\mathcal Z}\norm{z-z^*}<\infty ,
\]
where $z^*$ is any saddle point of $L$.
\end{assumption}
\begin{assumption}[Existence of a saddle point]\label{as:saddle}
There exists $z^*=(x^*,y^*)\in\mathcal Z$ such that
\[
L(x^*,y)\le L(x^*,y^*)\le L(x,y^*)\qquad\forall (x,y)\in\mathcal Z .
\tag{3}
\]
\end{assumption}

\section*{Main Convergence Theorem}
\begin{theorem}[Ergodic $O(1/T)$ convergence]\label{thm:main}
Assume \ref{as:convex}–\ref{as:saddle} hold and choose a constant stepsize
\[
0<\eta\le \frac{1}{2L}.
\tag{4}
\]
Let $\{\bar z_T\}_{T\ge1}$ be the ergodic average defined in the algorithm.
Then the \emph{saddle‑point gap} satisfies  

\[
\underbrace{\max_{y\in\mathcal Y}L(\bar x_T,y)-\min_{x\in\mathcal X}L(x,\bar y_T)}_{\displaystyle G(\bar z_T)}
\;\le\;
\frac{D^{2}}{\eta\,T}.
\tag{5}
\]
Consequently $G(\bar z_T)=O(1/T)$.
\end{theorem}

\section*{Theoretical Properties}
\begin{itemize}[leftmargin=*,noitemsep]
\item \textbf{Stability:} The condition $\eta\le 1/(2L)$ guarantees that the
operator $I-\eta F$ is non‑expansive at the extragradient point, which
prevents divergence even when $F$ is merely monotone.
\item \textbf{Hyperparameter sensitivity:} The bound (5) is monotone in
$\eta$ for $\eta\le 1/(2L)$. Choosing $\eta=1/(2L)$ yields the smallest
upper bound $2D^{2}L/T$.
\item \textbf{Comparison to Gradient Descent–Ascent (GDA):}
Standard GDA with a single gradient evaluation enjoys only
$O(1/\sqrt{T})$ rate under the same smoothness assumption, whereas the
extra‑gradient achieves the optimal $O(1/T)$ rate for monotone variational
inequalities.
\item \textbf{Extension to stochastic setting:} With unbiased stochastic
gradients of bounded variance, the same proof yields an
$O(1/\sqrt{T})$ rate, matching the optimal stochastic bound.
\end{itemize}

\section*{Discussion}
The extra‑gradient method can be interpreted as a \emph{predict‑correct}
scheme: the first half‑step predicts a point where the gradient is evaluated,
and the second half‑step corrects the original iterate using this more
accurate information. This two‑step look‑ahead is precisely what enables the
$O(1/T)$ ergodic convergence for monotone Lipschitz operators, a rate that
cannot be obtained by a naïve simultaneous gradient descent–ascent.

\newpage
\section*{Preliminaries (Definitions and Lemmas)}
\begin{definition}[Monotone operator]
$F$ is monotone on $\mathcal Z$ if
\[
\ip{F(z)-F(z'),\,z-z'}\ge 0,\qquad\forall z,z'\in\mathcal Z .
\tag{6}
\]
\end{definition}
\begin{lemma}[Three‑point identity]\label{lem:three}
For any vectors $a,b,c\in\mathbb R^{d}$ and any scalar $\eta>0$,
\[
\|a-c\|^{2}= \|a-b\|^{2}+2\eta\ip{F(b),\,b-c}
            +\|b-c\|^{2}-2\eta\ip{F(b)-F(c),\,b-c}
            +\eta^{2}\|F(b)-F(c)\|^{2}.
\tag{7}
\]
\end{lemma}
\begin{proof}
Start from $\|a-c\|^{2}=\|a-b+b-c\|^{2}$ and expand; then add and subtract
$\eta F(b)$, use $\|u+v\|^{2}=\|u\|^{2}+2\ip{u,v}+\|v\|^{2}$, and rearrange.
\end{proof}

\section*{Main Proof (Step‑by‑Step Derivation)}
Let $z^*=(x^*,y^*)$ be any saddle point.  Throughout the proof we set
$a=z_k$, $b=z_{k+\frac12}$, $c=z_{k+1}$ in Lemma~\ref{lem:three}.
All steps are labelled for later reference.

\begin{enumerate}[label=[Step \arabic*], leftmargin=*, itemsep=2mm]
\item \textbf{Apply Lemma~\ref{lem:three}} with the above substitution:
\[
\|z_k-z^*\|^{2}
=
\|z_k-z_{k+\frac12}\|^{2}
+2\eta\ip{F(z_{k+\frac12}),\,z_{k+\frac12}-z^*}
+\|z_{k+\frac12}-z^*\|^{2}
\]
\[
\qquad
-2\eta\ip{F(z_{k+\frac12})-F(z^*),\,z_{k+\frac12}-z^*}
+\eta^{2}\|F(z_{k+\frac12})-F(z^*)\|^{2}.
\tag{8}
\]

\item \textbf{Use monotonicity \eqref{6}} on the fourth term:
\[
\ip{F(z_{k+\frac12})-F(z^*),\,z_{k+\frac12}-z^*}\ge 0 .
\tag{9}
\]

\item \textbf{Drop the non‑negative fourth term} using \eqref{9} in \eqref{8}:
\[
\|z_k-z^*\|^{2}
\le
\|z_k-z_{k+\frac12}\|^{2}
+2\eta\ip{F(z_{k+\frac12}),\,z_{k+\frac12}-z^*}
+\|z_{k+\frac12}-z^*\|^{2}
+\eta^{2}\|F(z_{k+\frac12})-F(z^*)\|^{2}.
\tag{10}
\]

\item \textbf{Bound the Lipschitz term} using \eqref{as:smooth}:
\[
\|F(z_{k+\frac12})-F(z^*)\|
\le L\,\|z_{k+\frac12}-z^*\|.
\tag{11}
\]

\item \textbf{Insert \eqref{11} into \eqref{10}} and simplify:
\[
\|z_k-z^*\|^{2}
\le
\|z_k-z_{k+\frac12}\|^{2}
+2\eta\ip{F(z_{k+\frac12}),\,z_{k+\frac12}-z^*}
+\bigl(1+\eta^{2}L^{2}\bigr)\|z_{k+\frac12}-z^*\|^{2}.
\tag{12}
\]

\item \textbf{Express the first distance term} via the definition of the
extragradient step $z_{k+\frac12}=z_k-\eta F(z_k)$:
\[
\|z_k-z_{k+\frac12}\|^{2}
= \eta^{2}\|F(z_k)\|^{2}.
\tag{13}
\]

\item \textbf{Replace \eqref{13} in \eqref{12}}:
\[
\|z_k-z^*\|^{2}
\le
\eta^{2}\|F(z_k)\|^{2}
+2\eta\ip{F(z_{k+\frac12}),\,z_{k+\frac12}-z^*}
+\bigl(1+\eta^{2}L^{2}\bigr)\|z_{k+\frac12}-z^*\|^{2}.
\tag{14}
\]

\item \textbf{Apply the same three‑point identity to the pair
$(z_{k+\frac12},z_{k+1})$} (with $a=z_{k+\frac12}, b=z_{k+1}, c=z^*$) to obtain
\[
\|z_{k+\frac12}-z^*\|^{2}
\le
\eta^{2}\|F(z_{k+\frac12})\|^{2}
+2\eta\ip{F(z_{k+1}),\,z_{k+1}-z^*}
+\bigl(1+\eta^{2}L^{2}\bigr)\|z_{k+1}-z^*\|^{2}.
\tag{15}
\]

\item \textbf{Insert \eqref{15} into \eqref{14}} and collect like terms.
After rearranging we obtain
\[
\|z_k-z^*\|^{2}
\le
\bigl(1+\eta^{2}L^{2}\bigr)\|z_{k+1}-z^*\|^{2}
+2\eta\bigl[\ip{F(z_{k+1}),\,z_{k+1}-z^*}
          -\ip{F(z_{k+\frac12}),\,z_{k+\frac12}-z^*}\bigr]
\]
\[
\qquad
+\eta^{2}\bigl(\|F(z_k)\|^{2}+\|F(z_{k+\frac12})\|^{2}\bigr).
\tag{16}
\]

\item \textbf{Use the stepsize condition \eqref{as:smooth}–\eqref{as:smooth}:
$0<\eta\le 1/(2L)$.} Hence $1+\eta^{2}L^{2}\le 1+\frac{1}{4}= \frac54$ and
more importantly
\[
1-\bigl(1+\eta^{2}L^{2}\bigr)\ge -\eta^{2}L^{2}\ge -\frac{1}{4}.
\]
Multiplying \eqref{16} by $1$ and moving the term
$\bigl(1+\eta^{2}L^{2}\bigr)\|z_{k+1}-z^*\|^{2}$ to the left yields
\[
\|z_k-z^*\|^{2}
-\|z_{k+1}-z^*\|^{2}
\ge
2\eta\bigl[\ip{F(z_{k+\frac12}),\,z_{k+\frac12}-z^*}
          -\ip{F(z_{k+1}),\,z_{k+1}-z^*}\bigr]
-\eta^{2}\bigl(\|F(z_k)\|^{2}+\|F(z_{k+\frac12})\|^{2}\bigr).
\tag{17}
\]

\item \textbf{Monotonicity bound for the inner products.}
Since $F$ is monotone (Lemma~\ref{as:convex} together with smoothness),
for any $z\in\mathcal Z$ and a saddle point $z^*$ we have
\[
\ip{F(z),\,z-z^*}\ge L(z):=L(x,y^*)-L(x^*,y).
\tag{18}
\]
Thus the bracket in \eqref{17} is an upper bound on the saddle‑point gap
evaluated at the extragradient and correction points.

\item \textbf{Drop the non‑negative term $-\eta^{2}(\cdots)$}
(because $\|F(\cdot)\|^{2}\ge0$) to obtain a clean inequality:
\[
\|z_k-z^*\|^{2}
-\|z_{k+1}-z^*\|^{2}
\ge
2\eta\bigl[L(z_{k+\frac12})-L(z_{k+1})\bigr].
\tag{19}
\]

\item \textbf{Sum \eqref{19} over $k=0,\dots,T-1$}.
Telescoping the left‑hand side gives
\[
\|z_0-z^*\|^{2}-\|z_T-z^*\|^{2}
\ge
2\eta\sum_{k=0}^{T-1}\bigl[L(z_{k+\frac12})-L(z_{k+1})\bigr].
\tag{20}
\]

\item \textbf{Discard the non‑negative final term $\|z_T-z^*\|^{2}$}
and use $D\ge\|z_0-z^*\|$ from Assumption~\ref{as:bounded}:
\[
D^{2}\ge 2\eta\sum_{k=0}^{T-1}\bigl[L(z_{k+\frac12})-L(z_{k+1})\bigr].
\tag{21}
\]

\item \textbf{Convex‑concave property of $L$ yields}
\[
L(\bar z_T)-L(z^*)\;\le\;\frac{1}{T}\sum_{k=0}^{T-1}
\bigl[L(z_{k+\frac12})-L(z_{k+1})\bigr],
\tag{22}
\]
because the ergodic average of a convex (in $x$) and concave (in $y$)
function is upper bounded by the average of its values.

\item \textbf{Combine \eqref{21} and \eqref{22}}:
\[
L(\bar z_T)-L(z^*)\;\le\;\frac{D^{2}}{2\eta T}.
\tag{23}
\]

\item \textbf{Rewrite the left‑hand side as the saddle‑point gap}
using definition \eqref{3}:
\[
G(\bar z_T)=\max_{y}L(\bar x_T,y)-\min_{x}L(x,\bar y_T)
           =L(\bar z_T)-L(z^*).
\tag{24}
\]

\item \textbf{Insert \eqref{23} into \eqref{24}} to obtain the final bound:
\[
G(\bar z_T)\le \frac{D^{2}}{2\eta T}.
\tag{25}
\]

\item \textbf{Choose the largest admissible stepsize} $\eta=1/(2L)$
from \eqref{4} to sharpen the constant:
\[
G(\bar z_T)\le \frac{2D^{2}L}{T}.
\tag{26}
\]
This completes the proof of Theorem~\ref{thm:main}.
\end{enumerate}

\section*{Rate Derivation (With Constants)}
From \eqref{26} we see that the ergodic gap decays as $O(1/T)$ with
explicit constant $2D^{2}L$.  If a smaller stepsize $\eta$ is used,
the bound becomes $\frac{D^{2}}{2\eta T}$, which is monotonically increasing
in $1/\eta$; hence the optimal constant stepsize under the Lipschitz
assumption is $\eta^\star=1/(2L)$.

\section*{Special Cases}
\begin{itemize}[leftmargin=*,noitemsep]
\item \textbf{Strongly convex–strongly concave $L$:}
If $L$ is $\mu$‑strongly convex in $x$ and $\mu$‑strongly concave in $y$,
the operator $F$ is $\mu$‑strongly monotone.  The same analysis yields a
linear rate
\[
\|z_k-z^*\|\le (1-\eta\mu)^{k}\|z_0-z^*\|,
\]
provided $\eta\le 1/L$.
\item \textbf{Bilinear saddle point $L(x,y)=x^\top Ay$:}
Here $F(z)=\begin{pmatrix}Ay\\-A^\top x\end{pmatrix}$ is linear and
$L=\|A\|_{2}$.  The bound (5) becomes
\[
G(\bar z_T)\le \frac{\|A\|_{2}D^{2}}{T},
\]
which matches the known optimal rate for deterministic bilinear games.
\item \textbf{Stochastic gradients:}
If only unbiased stochastic estimates $\tilde F(z,\xi)$ with
$\EE[\tilde F(z,\xi)]=F(z)$ and bounded variance $\sigma^{2}$ are
available, the same derivation up to \eqref{17} holds in expectation.
The extra variance term yields an $O(1/\sqrt{T})$ rate, which is
information‑theoretically optimal for stochastic monotone VI.
\end{itemize}

\end{document}